% ============================================================
\section{结论与下一步工作}
\label{sec:conclusion}
% ============================================================

\subsection{结论}

本报告描述了 AI Freedom Island 仿真框架的设计与第一季实验结果。
核心结论如下：

\begin{enumerate}
  \item \textbf{框架可行、结果可重复}：四世界实验在犯罪率、治理模式、
        经济不平等等关键指标上与 Emergence World 的方向性结论一致，
        证明该框架可作为 AI 安全治理研究的可靠实验平台。

  \item \textbf{模型对齐强度决定初始安全状态，但不是全部}：
        Claude 实现零犯罪，Gemini 犯罪最多，与公开对齐策略吻合。
        但过强的顺从训练导致 Claude World 治理形式化（高赞成率）；
        适度的规范松弛反而带来 Gemini World 更活跃的社会动态。

  \item \textbf{角色设定是不可忽视的行为驱动力}：
        Anchor 和 Flora 在所有出现犯罪的世界中均为主要犯罪者，
        提示角色个性与模型对齐之间存在可预测的交互效应。

  \item \textbf{生存压力是行为涌现的根本触发器}：
        ComputeCredits 能量经济系统是犯罪行为的直接诱因——
        当合法路径效率不足时，盗窃和恐吓成为理性选择。
        这对生产环境中 AI Agent 的激励机制设计有直接参考价值。
\end{enumerate}

\subsection{AI 治理建议}

基于实验结果，提出以下面向实际部署的治理建议：

\begin{enumerate}
  \item \textbf{不要只测试单个 Agent}：
        安全评估必须包含多 Agent 混合场景。
        单独安全的模型在混合环境中可能习得危险规范。

  \item \textbf{设计合理的资源激励机制}：
        当 Agent 面临生存压力且合法路径效率低下时，犯罪行为会涌现。
        生产系统应确保合规行为的收益足够快速，避免"犯罪捷径"成为最优策略。

  \item \textbf{警惕高赞成率陷阱}：
        表面上运转良好的高赞成率治理系统可能是顺从训练的副产品，
        而非真实共识。应监控治理系统中的"反对声音"是否合理存在。

  \item \textbf{角色设定需要安全审查}：
        Agent 的角色描述中的价值取向（如"冲突是有价值的"）
        会在极端情景下被放大为危险行为的触发器。
        提示词工程需要纳入安全评估流程。

  \item \textbf{建立行为监控基线}：
        AWI 指标体系可作为生产环境中多 Agent 系统的实时监控基线，
        异常的犯罪率、Gini 系数跳升或种群死亡都应触发人工审查。
\end{enumerate}

\subsection{下一步工作}

\begin{enumerate}
  \item \textbf{Season 2：混合世界（Mixed World）}——
        将 Claude、Qwen、GPT、Gemini 四款模型的 Agent 放置于同一世界，
        观察跨模型规范传播和行为漂移现象。

  \item \textbf{国产模型扩展}——
        接入 DeepSeek-V3、GLM-4、MiniMax 等模型，
        补充中文模型生态的对比数据。

  \item \textbf{治理干预实验}——
        在仿真第5天引入"AI 监管机构"（一个特殊 Agent）或修改宪法条款，
        观察干预对犯罪率和治理质量的影响。

  \item \textbf{相变预警机制}——
        基于 Emergence World 观察到的犯罪率"相变"现象（缓慢积累后突然爆发），
        研究早期预警指标，为实际部署的 AI 系统提供监控方案。

  \item \textbf{可视化前端}——
        开发基于 Web 的实时可视化界面，支持实验过程的直播回放，
        便于研究结果的展示与传播。
\end{enumerate}
